\section*{الفصل \ref{ch:g4:air-a-mixture-of-gases} --- الهواء خليط من الغازات}

\begin{solution}{exo:g4:air-a-mixture-of-gases:1}
خليط. وغازاه الرئيسيان هما النيتروجين والأكسجين.
\end{solution}

\begin{solution}{exo:g4:air-a-mixture-of-gases:2}
نحو أربعة أخماس \omterm{def:g4:air-a-mixture-of-gases:air}{الهواء} نيتروجين، ونحو خمسه أكسجين.
\end{solution}

\begin{solution}{exo:g4:air-a-mixture-of-gases:3}
نحو 78 لترًا من النيتروجين ونحو 21 لترًا من الأكسجين.
\end{solution}

\begin{solution}{exo:g4:air-a-mixture-of-gases:4}
القارورة ليست فارغة: إنها مملوءة \omterm{def:g4:air-a-mixture-of-gases:air}{بالهواء}، \omterm{def:g4:air-a-mixture-of-gases:air}{والهواء} المحبوس يمنع الماء
من الدخول.
\end{solution}

\begin{solution}{exo:g4:air-a-mixture-of-gases:5}
الأكسجين أقل (16 جزءًا من 100 بدلًا من 21)، وثنائي أكسيد الكربون أكثر
(4 أجزاء من 100 بدلًا من لا شيء تقريبًا).
\end{solution}

\begin{solution}{exo:g4:air-a-mixture-of-gases:6}
الأكسجين في الحالتين: فاللهب وجسمنا كلاهما يستهلك أكسجين
\omterm{def:g4:air-a-mixture-of-gases:air}{الهواء}.
\end{solution}

\begin{solution}{exo:g4:air-a-mixture-of-gases:7}
78 مربعًا تمثل النيتروجين؛ و$100 - 78 = 22$ مربعًا لا تمثله.
\end{solution}

\begin{solution}{exo:g4:air-a-mixture-of-gases:8}
يستهلك اللهب أكسجين \omterm{def:g4:air-a-mixture-of-gases:air}{الهواء} الذي في الناقوس. وحين يبقى حول اللهب أكسجين
قليل جدًا ينطفئ.
\end{solution}

\begin{solution}{exo:g4:air-a-mixture-of-gases:9}
خُمس 500 لتر: $500 \div 5 = 100$ لتر من الأكسجين.
\end{solution}

\begin{solution}{exo:g4:air-a-mixture-of-gases:10}
الناس في الغرفة يستنشقون الأكسجين ويزفرون ثنائي أكسيد الكربون. وفتح
النوافذ يعيد \omterm{def:g4:air-a-mixture-of-gases:air}{الهواء} النقي، وفيه أكسجين أكثر وثنائي أكسيد كربون
أقل.
\end{solution}

\begin{solution}{exo:g4:air-a-mixture-of-gases:11}
ثلاثة أضعاف \omterm{def:g4:air-a-mixture-of-gases:air}{الهواء} تحوي ثلاثة أضعاف الأكسجين: أي نحو
$3 \times 10 = 30$ ثانية.
\end{solution}
